\section{Stable lengths and centralizer characters}\label{sec:lengths}

Our objective is a pairing of stable lengths that is additive on an
entire centralizer, even when its elements do not commute with one
another.  The geometric facts used here are standard: attainment of minimum
displacement under a geometric action and the structure of its
minimizing set are treated in Bridson--Haefliger
\cite[II.6.8(5) and II.6.10(2)]{BridsonHaefliger1999}; the Euclidean structure
on commuting translations is central to the Flat Torus Theorem
\cite[II.7.1]{BridsonHaefliger1999}.  We derive precisely the length
identities needed for the braid calculation directly from triangle
comparison, without using a splitting theorem.

\begin{lemma}\label{lem:comparison}
In a CAT(0) space, geodesic segments are unique and depend continuously on
their endpoints.  If $\alpha,\beta:[0,1]\to X$ are geodesics parametrized
proportionally to length, then
\begin{equation}\label{eq:distance-convexity}
 d(\alpha(t),\beta(t))
 \le (1-t)d(\alpha(0),\beta(0))+t d(\alpha(1),\beta(1)).
\end{equation}
If $m$ is the midpoint of $[v,w]$, then
\begin{equation}\label{eq:midpoint-comparison}
 d(z,m)^2\le \frac{d(z,v)^2+d(z,w)^2}{2}
                    -\frac{d(v,w)^2}{4}.
\end{equation}
For any four points $p,q,r,s$,
\begin{equation}\label{eq:quadrilateral}
 d(p,r)^2+d(q,s)^2
 \le d(p,q)^2+d(q,r)^2+d(r,s)^2+d(s,p)^2.
\end{equation}
\end{lemma}

\begin{proof}
For uniqueness, split one of two segments with the same endpoints at an
arbitrary point and use the other segment as the third side of the resulting
triangle.  Its comparison triangle is a straight segment, so comparison
identifies the two points with the same proportional parameter.  For two
segments with a common initial point, comparison bounds their distance at
parameter $t$ by $t$ times the distance of their terminal points.  The
analogous statement holds for a common terminal point.  Inserting the
segment from $\alpha(0)$ to $\beta(1)$ between $\alpha$ and $\beta$, and
using the triangle inequality, proves \eqref{eq:distance-convexity}.
This bound also proves continuous dependence on the endpoints; dependence
on the parameter follows from constant-speed parametrization.

The Euclidean identity for a median, together with comparison, gives
\eqref{eq:midpoint-comparison}.  Let $m$ now be the midpoint of $[p,r]$.
The triangle inequality and $(A+B)^2\le 2A^2+2B^2$ give
\[
 d(q,s)^2\le 2d(q,m)^2+2d(s,m)^2.
\]
Apply \eqref{eq:midpoint-comparison} to each term on the right to obtain
\eqref{eq:quadrilateral}.
\end{proof}

For an isometry $g$ and a base point $o$, define its \emph{stable length} by
\begin{equation}\label{eq:stable-length}
 \ell(g)=\lim_{n\to\infty}\frac{d(o,g^n o)}{n}.
\end{equation}
Here the limit exists: the sequence $a_n=d(o,g^n o)$ is subadditive, and,
for fixed $k$ and $n=qk+r$ with $0\le r<k$,
$a_n\le q a_k+a_r$.  Consequently
$\limsup a_n/n\le a_k/k$ for every $k$, and the limit equals
$\inf_{k\ge1}a_k/k$.  Changing $o$ changes $a_n$ by at most twice the
distance between the two base points.  Thus the length is independent of
$o$, is invariant under conjugation, and satisfies
\begin{equation}\label{eq:length-powers}
 \ell(g^m)=|m|\ell(g)\qquad(m\in\mathbb Z).
\end{equation}
For negative $m$, use $d(o,g^{-n}o)=d(o,g^n o)$; for positive $m$, use the
subsequence indexed by multiples of $m$ in \eqref{eq:stable-length}.

\begin{lemma}\label{lem:semisimple}
Suppose $G$ acts geometrically on $X$.  Every $g\in G$ attains its minimum
displacement, and
\[
 \ell(g)=\min_{x\in X}d(x,gx).
\]
If this number is positive, $g$ preserves a geodesic line and translates
along it by $\ell(g)$.  Moreover, $\ell(g)=0$ if and only if $g$ has finite
order.
\end{lemma}

\begin{proof}
Set $D=\inf_xd(x,gx)$, and choose $x_n$ with
$d(x_n,gx_n)\to D$.  Write $x_n=h_n p_n$ with $p_n$ in a compact set $K$
whose $G$-translates cover $X$.  After discarding finitely many terms,
$q_n=h_n^{-1}gh_n$ moves $p_n$ by at most $D+1$.  The closed
$(D+1)$-neighborhood $K'$ of $K$ is closed and bounded, hence compact.
Both $p_n$ and $q_np_n$ belong to $K'$, so $q_nK'\cap K'\ne\varnothing$.
Action properness gives only finitely many possibilities for $q_n$.
Pass to a subsequence on which $q_n=q$ and $p_n\to p\in K$.  Then
$d(p,qp)=D$.  A fixed one of the conjugators $h_n$ transfers $p$ to a
minimizer for $g$.

Let $x$ be such a minimizer.  If $D>0$, let $m$ be the midpoint of
$[x,gx]$; then $gm$ is the midpoint of $[gx,g^2x]$.  Comparison in the
triangle with vertices $x,gx,g^2x$ gives
\[
 D\le d(m,gm)\le \tfrac12 d(x,g^2x)\le D.
\]
Thus $d(x,g^2x)=2D$.  We verify that all consecutive orbit segments form
a line.  Suppose the concatenation from $x$ to $g^i x$ is geodesic.
In the triangle with vertices $x,g^i x,g^{i+1}x$, the points at distance
$D$ from $g^i x$ along the two incident sides are $g^{i-1}x$ and
$g^{i+1}x$.  Their distance is $2D$, by the preceding equality and
translation by $g^{i-1}$.  Their Euclidean comparison points can be at
distance $2D$ only if the comparison angle at $g^i x$ is $\pi$.
Hence $d(x,g^{i+1}x)=(i+1)D$.  Induction, followed by translation by
all powers of $g$, proves that the segments $[g^i x,g^{i+1}x]$,
$i\in\mathbb Z$, form an isometrically parametrized line.  In particular
$\ell(g)=D$.

If $D=0$, the minimizer is fixed, so $\ell(g)=0$.  Its stabilizer is finite
by action properness applied to the compact singleton $\{x\}$; thus $g$
has finite order.  Conversely, finite order immediately gives zero
stable length.
\end{proof}

For commuting isometries of a complete CAT(0) space, Rodenhausen
\cite[Theorem~14.6]{Rodenhausen2013} proves a parallelogram identity
for infimal displacements.  The next proof works directly with stable
lengths.

\begin{lemma}\label{lem:commuting-form}
For pairwise commuting isometries $g_1,\ldots,g_k$ of a CAT(0) space,
the function
\[
 v\longmapsto \ell(g^v),\qquad
 g^v=g_1^{v_1}\cdots g_k^{v_k},\qquad v\in\mathbb Z^k,
\]
extends uniquely to a continuous seminorm on $\mathbb R^k$ whose square
is a positive semidefinite quadratic form.  For commuting $g,h$, put
\begin{equation}\label{eq:pairing-definition}
 \langle g,h\rangle
   =\frac{\ell(gh)^2-\ell(g)^2-\ell(h)^2}{2}.
\end{equation}
These pairings are the bilinear form on every commuting family; in
particular, $\langle g^i,h^j\rangle=ij\langle g,h\rangle$.  Lengths and
pairings are invariant under simultaneous conjugation.
\end{lemma}

\begin{proof}
Write $\nu(v)=\ell(g^v)$ for integer $v$.  Commutation, the triangle
inequality, and \eqref{eq:length-powers} give
\begin{gather*}
 \nu(u+v)\le\nu(u)+\nu(v),\qquad
 \nu(mv)=|m|\nu(v),\\
 |\nu(u)-\nu(v)|\le\sum_{i=1}^k|u_i-v_i|\ell(g_i).
\end{gather*}
Define $\nu(v/q)=\nu(v)/q$ for integers $q>0$.  Homogeneity makes this
well-defined on $\mathbb Q^k$, and the displayed Lipschitz bound gives a
unique continuous extension to $\mathbb R^k$.  It is a seminorm.

Apply \eqref{eq:quadrilateral} to the four points
\[
 o,\quad g^{nu}o,\quad g^{n(u+v)}o,\quad g^{nv}o
\]
in that cyclic order.  Divide by $n^2$ and pass to the limit to obtain,
with $Q=\nu^2$,
\[
 Q(u+v)+Q(u-v)\le 2Q(u)+2Q(v).
\]
This initially holds for integer vectors and extends by scaling and
continuity to real vectors.  Apply it to $u+v,u-v$ and use
$Q(2u)=4Q(u)$ to obtain the opposite inequality.  Hence
\begin{equation}\label{eq:parallelogram}
 Q(u+v)+Q(u-v)=2Q(u)+2Q(v).
\end{equation}

For completeness, polarization here requires only elementary algebra.
Set $B(u,v)=(Q(u+v)-Q(u-v))/4$.  It is symmetric and continuous, and
\eqref{eq:parallelogram} gives
\[
 B(u+w,v)+B(u-w,v)=2B(u,v).
\]
Taking $w=u$ gives $B(2u,v)=2B(u,v)$, since $B(0,v)=0$.
Next take $u=(a+b)/2$, $w=(a-b)/2$ to obtain additivity in the first
variable.  Integer and rational homogeneity follow, and continuity gives
real homogeneity.  Symmetry gives the same properties in the second
variable.  Finally $B(u,u)=Q(u)\ge0$, so $B$ is positive semidefinite;
expanding $Q(u+v)$ recovers \eqref{eq:pairing-definition}.  Conjugation
invariance follows from that of stable length.
\end{proof}

The following strengthening of bilinearity is crucial.  The elements
$h_1,h_2$ in its statement need not commute with one another.
Real translation characters also underlie Bridson's obstruction to
proper semisimple actions on two-dimensional CAT(0) spaces
\cite[Proposition~3.1]{Bridson2001}.

\begin{lemma}\label{lem:centralizer-character}
For a geometric action and a fixed $g\in G$, the function
\[
 C_G(g)\longrightarrow\mathbb R,\qquad h\longmapsto\langle g,h\rangle
\]
is a group homomorphism.  Thus
$\langle g,h_1h_2\rangle=\langle g,h_1\rangle+\langle g,h_2\rangle$
whenever both $h_1$ and $h_2$ commute with $g$.
\end{lemma}

\begin{proof}
If $\ell(g)=0$, positive semidefiniteness in
Lemma~\ref{lem:commuting-form} implies $\langle g,h\rangle=0$ for every
$h\in C_G(g)$: the nonnegative quadratic
$\ell(g^t h)^2=2t\langle g,h\rangle+\ell(h)^2$, for integer $t$ of
both signs, forces the coefficient to vanish.

Now let $D=\ell(g)>0$ and choose a minimizer $x$.  By
Lemma~\ref{lem:semisimple}, $d(x,g^n x)=nD$.  Define
\[
 b(z)=\lim_{n\to\infty}\bigl(d(z,g^n x)-nD\bigr).
\]
The terms decrease by the triangle inequality and are bounded below by
$-d(z,x)$, so the limit exists.  The function $b$ is $1$-Lipschitz.

Fix $h\in C_G(g)$ and $z,w\in X$.  Put
$p_n=g^n x$ and $q_n=h^{-1}g^n x$.  Since $h$ commutes with $g$,
$d(p_n,q_n)=d(x,h^{-1}x)$ is constant.  Two orderings in
\eqref{eq:quadrilateral} give
\begin{align*}
 \bigl|d(z,p_n)^2+d(w,q_n)^2-d(w,p_n)^2-d(z,q_n)^2\bigr|
 \le d(z,w)^2+d(p_n,q_n)^2.
\end{align*}
Every long distance on the left is $nD+O(1)$.  More precisely,
$d(z,p_n)=nD+b(z)+o(1)$ and
$d(z,q_n)=d(hz,p_n)=nD+b(hz)+o(1)$, and similarly for $w$.
Dividing the inequality by $2nD$ and passing to the limit therefore gives
\begin{equation}\label{eq:busemann-difference}
 b(hz)-b(hw)=b(z)-b(w).
\end{equation}
Consequently $\chi(h)=b(x)-b(hx)$ is a homomorphism on $C_G(g)$:
apply \eqref{eq:busemann-difference} with $z=kx,w=x$ to compute
$\chi(hk)=\chi(h)+\chi(k)$.  Also $b(x)=0$, $b(gx)=-D$, and hence
$\chi(g)=D$.  Lipschitz continuity, applied to $h^n$, gives
\[
 n|\chi(h)|=|b(x)-b(h^n x)|\le d(x,h^n x),
 \qquad\text{so } |\chi(h)|\le\ell(h).
\]
Apply this last inequality to $g^t h$ for every integer $t$.  By
Lemma~\ref{lem:commuting-form},
\[
 (tD+\chi(h))^2
 \le t^2D^2+2t\langle g,h\rangle+\ell(h)^2.
\]
After cancelling the quadratic terms, the inequality for arbitrarily
large positive and negative $t$ forces
$D\chi(h)=\langle g,h\rangle$.  Thus the pairing is a multiple of the
homomorphism $\chi$, as asserted.
\end{proof}
